% Lösungen Trigonometrie · Teildreiecke und Vermessung · Prüfungs-Fokus MSA (zusammenbau v0.6)
\einheitenkopf{Einheit 3 · Rechtwinklige Teildreiecke in Figuren und Vermessung}

\erg{1}{a) Teildreieck aus Höhe, linkem Schenkel und linkem Überstand: Schenkel H, Höhe G, Überstand A \quad b) Teildreieck: Schenkel H, Höhe G. $\mathrm{sin}\,63^\circ = \frac{h}{7}$, $h = 7 \cdot \mathrm{sin}\,63^\circ \approx 6{,}2$ cm \quad c) $s = 4 : \mathrm{sin}\,62^\circ \approx 4{,}5$ m}
\erg{2}{a) $96{,}5 \cdot \mathrm{sin}\,27^\circ \approx 43{,}81$ m; Hügel $\approx 45{,}2$ m; Oberkante $\approx 54{,}2$ m \quad b) $148 \cdot \mathrm{sin}\,24^\circ \approx 60{,}20$ m; Hügel $\approx 61{,}9$ m; Oberkante $\approx 69{,}4$ m \quad c) Hilfsdreieck: Katheten $18 - 7 = 11$ (an $\alpha$) und $16$; $\mathrm{tan}\,\alpha = \frac{16}{11}$, $\alpha \approx 55{,}5^\circ$ \quad d) Hilfsdreieck: Katheten $13 - 5 = 8$ (an $\alpha$) und $10$; $\mathrm{tan}\,\alpha = \frac{10}{8}$, $\alpha \approx 51{,}3^\circ$ \quad e) $BF = 96 \cdot \mathrm{sin}\,41^\circ \approx 62{,}98$ cm, $AF = 96 \cdot \mathrm{cos}\,41^\circ \approx 72{,}45$ cm; im Dreieck $CBF$: $BC = BF : \mathrm{cos}\,57^\circ \approx 115{,}6$ cm \quad f) $BF = 86 \cdot \mathrm{sin}\,33^\circ \approx 46{,}84$ cm, $AF = 86 \cdot \mathrm{cos}\,33^\circ \approx 72{,}13$ cm; im Dreieck $CBF$: $BC = BF : \mathrm{cos}\,62^\circ \approx 99{,}8$ cm \quad g) Dreieck $AFC$, rechter Winkel bei $F$: $\mathrm{sin}\,19^\circ = \frac{h}{AC}$, $AC = 21{,}4 : \mathrm{sin}\,19^\circ \approx 65{,}7$ cm \quad h) Dreieck $AFC$, rechter Winkel bei $F$: $\mathrm{sin}\,23^\circ = \frac{h}{AC}$, $AC = 15{,}6 : \mathrm{sin}\,23^\circ \approx 39{,}9$ cm \quad i) $\angle CAD = 90^\circ - 59^\circ = 31^\circ$, $\angle DAB = 72^\circ - 31^\circ = 41^\circ$; $BD = 610 \cdot \mathrm{tan}\,41^\circ \approx 530$ m \quad j) $\angle CAD = 90^\circ - 63^\circ = 27^\circ$, $\angle DAB = 71^\circ - 27^\circ = 44^\circ$; $BD = 380 \cdot \mathrm{tan}\,44^\circ \approx 367$ m \quad k) Gebraucht: $AD$, $\delta$ und der rechte Winkel bei $E$. Reihenfolge: $AE = AD \cdot \mathrm{sin}\,\delta$, $DE = AD \cdot \mathrm{cos}\,\delta$, Flächeninhalt = die Hälfte von $AE \cdot DE$. \quad l) Gebraucht: $AD$, $\alpha$ und der rechte Winkel bei $F$. Reihenfolge: $DF = AD \cdot \mathrm{sin}\,\alpha$, $AF = AD \cdot \mathrm{cos}\,\alpha$, Flächeninhalt = die Hälfte von $AF \cdot DF$.}
